\appendix
\section{Worked geodesic calculations and possible replacements}
For trace $t=3$, the expanding eigenvalue is $(3+\sqrt5)/2$ and
\[
 N=\frac{7+3\sqrt5}{2},\qquad
 \ell=2\log\frac{3+\sqrt5}{2}.
\]
This is a legitimate hyperbolic length, but not $\log p$.
For $p=2$, the proposed length $\log2$ would require
$t=3/\sqrt2$, equivalently $2t^2=9$, which is impossible even modulo two.
The Lean companion verifies this last exact integer obstruction.
For arbitrary $p^k$, the integer argument in Theorem~\ref{thm:norm}
is stronger than a numerical search over traces.

A Hecke operator attached to determinant-$p$ double cosets is a different
object from a determinant-one deck transformation. A replacement must
specify the ambient space, coefficient system, trace class or regularized
trace, and the precise arithmetic distribution. Reusing the same symbol
for these objects does not construct an identification.
An adelic framework may encode prime norms intrinsically; its trace and
positivity obligations must then be proved in that framework.

\section{Character examples and induction bookkeeping}
The nonprincipal character modulo three has values
$\chi(1)=1$, $\chi(2)=-1$, $\chi(0)=0$ and is odd.
For scalar weight-zero automorphy under $\Gamma_0(3)$,
\[
 f(\gamma z)=\chi(d)f(z),
\]
the element $-I$ fixes $z$ and has $\chi(-1)=-1$. Hence every such
function would satisfy $f(z)=-f(z)$ and vanish. The proposed scalar bundle
on the projective quotient is therefore unavailable as written. One can
work on a different group or choose a compatible nonzero weight; either
change requires rebuilding the operator and trace identity.

Modulo five, the quadratic character has $\chi(2)=-1$.
Consequently the one-prime diagonal entry $\chi(2)2^{-1/2}$ is negative.
Its $1\times1$ minor is already negative, so an unrestricted assertion of
total nonnegativity of these transfer matrices is false. A positive state
constructed as $TT^*/\tr(TT^*)$ discards phase information in this diagonal
model: the nonzero character phases disappear from $TT^*$ at real $s$, although the ramification zero mask remains. Thus positivity
of that state cannot by itself recover character-sensitive Euler factors.

For $C_2\times C_2$, the three order-two subgroups each admit two linear
characters. All six induced characters have degree two; induction from
the trivial subgroup gives degree four. Negative integer coefficients do
not fix the parity obstruction. Rational coefficients can change parity,
which explains why rational and integral induction cannot be interchanged.
For a monomial representation $\operatorname{Ind}_H^G\theta$, the appropriate
Artin formalism yields a Hecke $L$-function over the fixed field of $H$;
it does not automatically reduce to a product of Dirichlet functions over
$\mathbb Q$. Every extension must retain the base field and infinity type.

\section{Local algebra and topology checks}
For the upper-Iwahori order
\[
 \mathcal O=\begin{pmatrix}R&R\\\pi R&R\end{pmatrix},\quad
 S=\begin{pmatrix}0&\pi\\1&0\end{pmatrix},
\]
$E_{12}\in\mathcal O$, whereas
$SE_{12}S^{-1}=\pi^{-1}E_{21}\notin\mathcal O$.
The displayed shift orientation therefore does not normalize this order.
Changing an orientation may repair this isolated calculation, but all
subsequent products and gluing maps must be rechecked with that convention.

For $R=\C[[t]]$ and $\lambda\in\C^\times$, there is no unital
$\C$-algebra homomorphism sending $t$ to $\lambda I$ in a nonzero matrix
algebra: $t-\lambda$ is a unit in $R$ and would be sent to zero.
Thus substituting a unit complex phase for a formal uniformizer is not a
representation of this allowed complete DVR. An algebra of Laurent
polynomials can admit unitary phase evaluation, but it is a different
algebra and needs its own construction.

A dessin on a sphere with two vertices and three edges has three faces,
by Euler's formula $V-E+F=2$, not two. Repeating a vertex along a face
boundary does not create a new primitive vertex idempotent. In particular,
a three-vertex path algebra with arrows between distinct vertices cannot
have the claimed surjection to a six-vertex radical-square-zero cycle
algebra: compose with the target's semisimple quotient $k^6$. Arrows map
to zero, since the orthogonal source idempotents force their endpoint
supports to be disjoint in the commutative target. The image is generated
by at most three orthogonal idempotents, so it cannot be all of $k^6$.
A valid replacement may require a cover or unfolded category rather than
a quotient of the stated algebra.

\section{Why weak convergence and boundary analogies are insufficient}
Let $P_n$ be projection onto the $n$th basis vector of $\ell^2$.
For every $x\in\ell^2$, $\|P_nx\|=|x_n|\to0$, so $P_n\to0$ strongly,
but $\|P_n\|_1=1$ for every $n$. Moreover
$\det(I-P_n)=0$ while $\det(I)=1$.
Thus even strong operator convergence does not imply determinant
convergence. Trace-norm control is a natural sufficient condition, and
must be established for the actual operator family, locally uniformly
in $s$, before invoking Fredholm determinant continuity.

For a Phragm\'en--Lindel\"of argument, one needs uniform control of
approximants as the approximation level changes. Polynomial growth of
each individual polynomial has constants and degrees depending on $N$;
it does not supply a common bound. Dividing by $(1+|s|^2)^M$ produces a
nonholomorphic expression, to which the maximum modulus principle does
not apply. Dividing by $(1+s^2)^M$ instead introduces poles at $s=\pm i$;
a denominator must be holomorphic and nonvanishing on the actual domain.
Even after selecting such a denominator, the necessary uniform boundary
and horizontal-side estimates are substantive mathematical work.

Likewise, a bound on a density of geodesics is not a bound on a weighted
analytic remainder. The full test-function dependence, prime powers,
scattering, elliptic and parabolic contributions must be retained.
On a cusped finite-area surface the infimum of the injectivity radius
is zero. It cannot provide the positive global mesh bound claimed in the
source completion argument.

\section{Spectral approximation can conceal the original hypothesis}
\begin{proposition}[Canonical-product converse for the Riemann target]
If RH holds, the entire order-one function $\xi$ admits locally uniform
approximations of the form in Theorem~\ref{thm:closure} with diagonal
self-adjoint $H_N$.
\end{proposition}
\begin{proof}
The Hadamard factorization of $\xi$, which has no zero at zero, is
\[
 \xi(s)=e^{a+bs}\prod_\rho(1-s/\rho)e^{s/\rho},
\]
with an enumeration for which the canonical product converges locally
uniformly. Under RH every $\rho=1/2+i\gamma$ has real $\gamma$.
Each finite product is an exponential of an affine function times a
constant times $\prod_{j\le N}(s-\rho_j)$; absorb the nonzero constant
into the exponential. Take $H_N=\operatorname{diag}(\gamma_1,\ldots,\gamma_N)$.
\end{proof}
Together with Theorem~\ref{thm:closure}, this shows why existence of the
approximants, without an independent arithmetic construction, may merely
restate RH. It is not progress to import the unknown zero locations into
the input spectrum and then prove that the constructed spectrum is real.
This proposition is a classical consequence, not a novel proof of RH.

\section{Hermitian dilation and the spectral parameter}
A common attempt to preserve arithmetic input while obtaining self-adjointness
is the dilation, for square $T(s)\in M_d(\C)$,
\[
 K(s)=\begin{pmatrix}0&T(s)\\T(s)^*&0\end{pmatrix}.
\]
It is Hermitian pointwise. Singular-value decomposition shows that its
spectrum is $\{\pm\sigma_j(T(s)):1\le j\le d\}$, with multiplicities.
Consequently
\[
 \det(\lambda I-K(s))=\prod_{j=1}^d(\lambda^2-\sigma_j(T(s))^2),
 \qquad \det(-K(s))=(-1)^d|\det T(s)|^2.
\]
These formulas retain singular values and erase determinant phase.
For a nonconstant holomorphic scalar $f$, $|f|^2$ cannot be holomorphic:
a real-valued holomorphic function is constant, and constant modulus then
forces $f$ constant. Thus the zero-matching scalar $|f|^2$ is not a
holomorphic replacement for $f$. Matrix nonholomorphy alone does not imply
that every scalar expression formed from it is nonholomorphic; each
proposed determinant identity must be examined separately.
For example, $T(s)=\operatorname{diag}(e^s,e^{-s})$ is nonconstant but its
dilation determinant is constantly one.

For the prime-diagonal model, the dilation eigenvalues are exactly
$\pm|\chi(p)|p^{-\Rea s}$. Both $\Im s$ and the nonzero character phases
disappear. This cannot encode the proposed arithmetic zero ordinates.
The source's parameter mismatch is visible where an $s$-dependent
transfer is converted to a Hermitian/modular Hamiltonian and then used as
though it had fixed eigenvalues in a characteristic polynomial
\cite{surface}, source lines 11465--11484 and 12200--12215; see also
lines 13636--13643 for the explicit skew-Hermitian-part extraction.
Its numerical code fixes a particular evaluation point instead, so the
rigidity proposition does not itself refute that distinct operation.
The earlier diagonal diagnostic obstruction addresses it directly.

A permitted architecture is a fixed self-adjoint $H_N$ with spectral
variable entering through $sI-(1/2+iH_N)$. Alternatively one may use a
holomorphic pencil that is Hermitian only on its real parameter axis,
but this condition does not force real spectral roots. Already the
scalar polynomial $h(t)=t^2+1$ is real for real $t$, whereas
$z-h(z)=0$ has the nonreal solutions $(1\pm i\sqrt3)/2$.
An infinite-dimensional operator model would additionally need explicit
domains, closure, adjoint and determinant regularization. No finite-matrix
rigidity statement here is asserted to settle variable-domain unbounded
operator constructions.

\section{Positive pencils, resolvent control, and a Herglotz test}
The positive generalized pencil is precisely a fixed self-adjoint
construction in different coordinates. Put $H=B^{-1/2}AB^{-1/2}=H^*$.
Then
\[
 A-zB=B^{1/2}(H-zI)B^{1/2},
\]
and therefore
\[
 \det(A+i(s-1/2)B)=i^d\det(B)
 \det\bigl(sI-(\tfrac12I+iH)\bigr).
\]
The positive metric $B$ must be supplied by the arithmetic construction;
its existence is not inferred from a desire for real roots.

There is also a quantitative test. If $B\ge bI$ with $b>0$ and $z$ is
nonreal, then for every $v$,
\[
 \|(A-zB)v\|\,\|v\|\ge |v^*(A-zB)v|
 \ge |\Im z|v^*Bv\ge b|\Im z|\,\|v\|^2.
\]
For $v\ne0$, divide by $\|v\|$. In finite dimension this proves
invertibility and the operator-norm bound
\[
 \|(A-zB)^{-1}\|\le\frac1{b|\Im z|}
 =\frac1{b|\Rea s-1/2|}.
\]
Uniform approximation arguments must track dimension, normalization and
$b_N$; a bound for each individual pencil does not by itself provide
uniform determinant convergence.

A nonlinear extension uses a matrix-valued Herglotz condition: if $M(z)$
is holomorphic upstairs and
\[
 \Im M(z):=\frac{M(z)-M(z)^*}{2i}>0,
\]
then $M(z)v=0$ would contradict $v^*\Im M(z)v>0$. Thus $\det M(z)\ne0$
there. A reflected lower-half-plane extension
$M(\bar z)=M(z)^*$ is likewise invertible. A quantitative lower bound
$\Im M(z)\ge\eta(z)I>0$ gives $\|M(z)^{-1}\|\le1/\eta(z)$ by the
same argument. For the positive linear pencil the Herglotz sign is
$M(z)=zB-A$, not $A-zB$.
Semidefinite positivity alone allows $M(z)=\operatorname{diag}(z,0)$,
which is always singular. Without reflection, $M(z)=z+i$ is strictly
Herglotz upstairs but has a lower-half-plane zero. Possible real-axis
poles and the arithmetic determinant identity remain separate obligations.

For a candidate normalized identity written in the $z$-coordinate,
$F_N(z)/g_N(z)=\det(A_N-zB_N)$ with a prescribed nowhere-zero $g_N$,
a necessary diagnostic, for $d_N>0$, is
\[
 -\partial_z\log(F_N/g_N)
 =\tr\bigl((A_N-zB_N)^{-1}B_N\bigr)
 =\sum_{j=1}^{d_N}\frac1{\lambda_{j,N}-z}.
\]
Its imaginary part is strictly positive for $\Im z>0$.
A certified violation rejects that particular determinant identity;
finitely many successful samples cannot prove it. One must not choose
$g_N$ using unknown zero locations or assume this Herglotz positivity as
a substitute for establishing the arithmetic bridge.

\section{Quantitative arithmetic approximation and prime regularization}
For the nonprincipal-character approximants, a compact set
$K\subset\{\Rea s>0\}$ has
$\delta=\min_K\Rea s>0$, $M=\max_K|s|$, and
$G=\max_K|(q/\pi)^{(s+a)/2}\Gamma((s+a)/2)|<\infty$.
The completed approximation error is at most
$GB(1+M/\delta)N^{-\delta}$ uniformly on $K$.
This is a fixed-character compact bound, not a uniform-in-conductor GRH estimate.
To verify the Abel formula's sign, first sum between integers $N<T$:
\[
 \sum_{N<n\le T}\chi(n)n^{-s}
 =A(T)T^{-s}-A(N)N^{-s}
   +s\int_N^T A(x)x^{-s-1}\,dx.
\]
Then $A(T)T^{-s}\to0$. Bounding $|A|$ by $B$ gives exactly the main-text estimate.
For zeta, $\lfloor x\rfloor=x-\{x\}$ in the same summation formula yields
\[
 \zeta(s)=\frac{s}{s-1}-s\int_1^\infty\{x\}x^{-s-1}\,dx.
\]
The finite sum formula gives the stated remainder after subtracting $Z_N$.
The integral and its derivatives converge uniformly on compact subsets
of the positive half-plane, since logarithmic powers are integrable
against $x^{-1-\delta}$. This justifies holomorphy and continuation.
Multiplying the remainder by $s(s-1)\pi^{-s/2}\Gamma(s/2)/2$
gives the compact-uniform $\xi_N$ error, including at $s=1$.
In contrast, multiplying a finite Euler product by $s(s-1)$ would force
a zero at $1$ for every truncation, although $\xi(1)=1/2$.

A second exact interface isolates what a prime-based repair would need.
For a nonprincipal primitive $\chi$, define on $\Rea s>1$
\[
 P_\chi(s)=\sum_p\chi(p)p^{-s},\qquad
 R_\chi(s)=\sum_p\sum_{k\ge2}\frac{\chi(p)^k}{k p^{ks}}.
\]
On each compact subset of $\Rea s>1/2$, the second series converges
absolutely and uniformly. Indeed, for $\Rea s\ge\delta>1/2$,
\[
 \sum_{k\ge2}\frac{|\chi(p)|^k}{k p^{k\delta}}
 \le\frac{p^{-2\delta}}{1-2^{-\delta}},
\]
and the prime sum is dominated by the convergent integer series
$\sum_{n\ge2}n^{-2\delta}$. Thus $R_\chi$ is holomorphic there.
On $\Rea s>1$, $L(s,\chi)=\exp(P_\chi(s)+R_\chi(s))$.

\begin{proposition}[The missing prime-linear term]\label{prop:prime-linear}
For a nonprincipal primitive character, $P_\chi$ has a single-valued
holomorphic continuation to $\Rea s>1/2$ if and only if
$L(s,\chi)$ is zero-free there.
\end{proposition}
\begin{proof}
If such a continuation exists, its exponential with $R_\chi$ agrees
with $L$ on $\Rea s>1$ and hence on the connected half-plane by the
identity theorem. The exponential is nonzero. Conversely, if $L$ is
nonzero, the simply connected half-plane admits a holomorphic logarithm
of $L$. Choose its additive constant to agree with the absolutely
convergent Euler logarithm on $\Rea s>1$. Subtract $R_\chi$ to obtain
the required continuation of $P_\chi$.
\end{proof}
This proposition isolates a precise analytic obligation; it does not prove it.
Regularizing away the terms $k\ge2$ supplies no proof that the missing
prime-linear function extends. The principal zeta case must first remove
the pole, for example by using $(s-1)\zeta(s)$; an unqualified logarithm
of $\zeta$ cannot be holomorphic through $1$.
The functional equation relates a fixed character's zeros across the
critical line after complex conjugation, so a full right-half-plane
nonvanishing theorem would give the desired Dirichlet GRH conclusion.

\section{Exact arithmetic behind the finite-level zero certificate}
This section gives all parameters needed to reproduce
Theorem~\ref{thm:finite-offline}. Write $C=|\Rea c|+|\Im c|<12$.
Use the following rational logarithm approximations $\ell_n$ and error
bounds $\delta_n$:
\[
\begin{array}{c|c|c}
 n&\ell_n&\delta_n\\\hline
 2&6931471805599453/10^{16}&10^{-15}\\
 4&13862943611198906/10^{16}&2\cdot10^{-15}\\
 5&16094379124341004/10^{16}&10^{-15}
\end{array}
\]
For $y=(n-1)/(n+1)$, set
\[
 L_n=2\sum_{k=0}^{59}\frac{y^{2k+1}}{2k+1},\qquad
 T_n=\frac{2y^{121}}{121(1-y^2)}.
\]
The positive series for $\log n$ gives $L_n\le\log n\le L_n+T_n$.
Exact rational comparisons give
$\ell_n-\delta_n\le L_n$ and $L_n+T_n\le\ell_n+\delta_n$.
They also give $0<\log2<1$ and $0<\log4,\log5<2$.

Define Gaussian rationals
\[
 E_n=\sum_{k=0}^{120}\frac{(-c\ell_n)^k}{k!},\qquad
 \tau=\frac{3^{24}24^{121}}{121!}<10^{-18},\qquad
 \eta_n=\tau+2C\delta_n.
\]
Since $|c\ell_n|\le C\ell_n<24$, the exponential series tail is
at most $e^{24}24^{121}/121!<\tau$. Here $e<3$ follows directly
from its series. Also
\[
 |e^{-c\log n}-e^{-c\ell_n}|
 \le e^{-\Rea(c)\ell_n}|c|\delta_n e^{|c|\delta_n}
 \le2C\delta_n,
\]
because $\ell_n,\Rea c>0$, $C\delta_n\le1/2$, and
$e^u\le(1-u)^{-1}\le2$ for $0\le u\le1/2$.
Thus $|n^{-c}-E_n|\le\eta_n$.

Put
\[
 F=1-E_2+E_4-E_5,\qquad
 D=\ell_2E_2-\ell_4E_4+\ell_5E_5.
\]
All entries are exact Gaussian rationals. With
$\|(a,b)\|_1=|a|+|b|$, the script checks the rational inequalities
\[
 \|F\|_1+\eta_2+\eta_4+\eta_5<10^{-6},\qquad
 |\Rea D|-\sum_{n\in\{2,4,5\}}(\delta_n+\ell_n\eta_n)>1.
\]
The first bounds $|f(c)|$. For the second, note that $|n^{-c}|\le1$
and
\[
 |(\log n)n^{-c}-\ell_nE_n|
 \le\delta_n+\ell_n\eta_n.
\]
This proves $|f'(c)|>1$. On the entire closed disk, $\Rea s>0$,
so
\[
 |f''(s)|\le(\log2)^2+(\log4)^2+(\log5)^2<1+4+4=9.
\]
Integrating the second derivative along the line segment from $c$ to $s$
justifies the Taylor remainder bound $9|s-c|^2/2$.
The exact strict boundary margin is
\[
 r-10^{-6}-\tfrac92r^2=\frac{19791}{200000000}>0.
\]
Consequently $|f(s)|$ is greater than that margin on $|s-c|=r$.
Any zero-free holomorphic positive-pencil model $G$ on a neighborhood of
this disk must have boundary error at least $\min_{|s-c|=r}|f(s)|$;
otherwise another application of Rouch\'e would contradict its zero count.
This is a quantitative rejection of a prescribed finite-level model,
not a proof about the zero locations of the limiting $L$-function.

The certificate program
\path{scripts/check_arithmetic_zero.py}
uses only Python integer and Fraction arithmetic.
A separate independent exact-arithmetic reproduction passed. Lean verifies
the disk's location and the last rational margin; the logarithm/exponential
error estimates and the Rouch\'e inference remain explicit prose proofs,
not claims of a complete kernel-checked zero-existence theorem.

\section{Jensen companions: exact normalization and the failed moment argument}
Our $\gamma_j$ equals one eighth of the convention in equation~1
of~\cite{jensen}, whose generating function is $8\xi(1/2+z)$.
This common positive factor changes no polynomial zeros. Entirety of
$F$ follows directly from the even Taylor series: its absolute series
at $|w|=R$ is the absolute series for $\xi(1/2+z)$ at $|z|=\sqrt R$.
The finite-product inequality in the main-text convergence proof is
also verified in Lean. The analytic summation and $\xi$ specialization
remain prose arguments, using classical properties of $\xi$.
For an explicit coefficient interface, the classical representation
in~\cite[equations 1.8--1.11]{romik} gives
\[
 \Phi(u)=2\sum_{n\ge1}
 (2\pi^2 n^4 e^{9u/2}-3\pi n^2e^{5u/2})e^{-\pi n^2e^{2u}},\qquad
 \gamma_j=\frac{j!}{(2j)!}\int_{\R}u^{2j}\Phi(u)\,du.
\]
The kernel is even. Each summand is positive for $u\ge0$ because
$2\pi n^2e^{2u}>3$, so the moments are strictly positive.
Its rapid decay permits differentiation of the Fourier integral.
This arithmetic kernel representation is an external classical input,
not a new Lean theorem in this project.

\subsection{The coefficient construction}
Write $p(x)=J_d(x)/\gamma_d=x^d+\sum_{i<d}a_ix^i$.
The companion $C$ has $C_{j+1,j}=1$ for $0\le j<d-1$,
last column $C_{i,d-1}=-a_i$, and all other entries zero.
For $1\le k\le d$ Newton's recurrence is
\[
 m_k+\sum_{i=1}^{k-1}a_{d-i}m_{k-i}+ka_{d-k}=0,
\]
and for $k>d$ it is
\[
 m_k+\sum_{i=1}^d a_{d-i}m_{k-i}=0.
\]
Together with $m_0=d$, this determines the entire required matrix
$H=(m_{i+j})_{0\le i,j<d}$. Multiplication in the quotient algebra
$\R[x]/(p)$ gives its bilinear form
\[
 h(u,v)=\tr(M_{uv}),\qquad h(xu,v)=h(u,xv),
\]
so $HC=C^TH$ without any zero-location assumption.
Here $M_{uv}$ means multiplication by $uv$ in that finite quotient.
When roots are real and distinct, the evaluation matrix $V$ gives
$H=V^TV>0$. Repeated roots make $V$ singular; the semidefinite
Hermite criterion still applies, but inverse square roots do not.
Nonnegative leading principal minors alone are not a PSD certificate;
all principal minors, or a valid alternative test, are required.

Assume now $H>0$. Then
\[
 S=H^{1/2}(-C)H^{-1/2}=S^T>0,\qquad
 K=\begin{pmatrix}0&\sqrt{dS}\\\sqrt{dS}&0\end{pmatrix}=K^T.
\]
Indeed the similarity is symmetric by $HC=C^TH$, and its eigenvalues
are negatives of the roots of $p$, all positive. The block determinant
and companion characteristic polynomial give
\begin{align*}
 \det(tI_{2d}-K)&=\det(t^2I_d-dS)
                  =\det(t^2I_d+dC),\\
 J_d(-t^2/d)&=\gamma_d(-1/d)^d\det(t^2I_d+dC).
\end{align*}
The identity holds also at $t=0$, by polynomial equality.
Using a companion for the scaled $Q_d$ instead would remove the extra
$d$ inside the square root. With our stated convention it cannot be omitted.
Thus the arithmetic data and approximation rate are specified; the
unproved universal positivity is precisely where the RH difficulty lies.
This construction supplies no general Dirichlet, Hecke, Dedekind or
Artin conclusion. Their symmetry and coefficient interfaces must be
established separately.

\subsection{Why ordinary moment positivity fails}
For $p=x^2+bx+a$, the companion and Newton-sum metric are
\[
 C=\begin{pmatrix}0&-a\\1&-b\end{pmatrix},\qquad
 H=\begin{pmatrix}2&-b\\-b&b^2-2a\end{pmatrix}.
\]
Direct calculation gives $HC=(HC)^T$, $\det H=b^2-4a$, and
\[
 u^THu=2(u_0-bu_1/2)^2+\tfrac12(b^2-4a)u_1^2.
\]
These identities and positivity under $b^2-4a>0$ are formalized in Lean.
For $J_2=\gamma_0+2\gamma_1x+\gamma_2x^2$, $\gamma_2>0$,
\[
 \det H=\frac{4(\gamma_1^2-\gamma_0\gamma_2)}{\gamma_2^2}.
\]
The unsmoothed atomic example has $\mu_0=1$, $\mu_2=\mu_4=1/10$.
Its ordinary even-moment Hankel matrix and Jensen metric are respectively
\[
 \begin{pmatrix}1&1/10\\1/10&1/10\end{pmatrix}>0,\qquad
 \begin{pmatrix}2&-12\\-12&-96\end{pmatrix},\quad\det H=-336.
\]
Gaussian smoothing gives the strictly positive densities in the main-text
counterexample. The standard Gaussian moments yield
$\mu_2=v+1/10$ and $\mu_4=3v^2+3v/5+1/10$.
With $\gamma_n=n!\mu_{2n}/(2n)!$, the negative gap is independent of $v$.
Lean checks that exact algebraic cancellation; the measure integral and
Fourier-transform identities remain prose. Thus even smooth positive
kernel information cannot replace arithmetic-specific Jensen inequalities.

The exact program \path{scripts/check_jensen_metric.py} verifies the
Newton recurrence, companion symmetry, positive LDL pivots for the
$\cosh\sqrt w$ model at degrees $1$ through $8$, and the negative and
singular witnesses. Its finite coefficient-bound checks supplement the
universal proof; they do not replace it. No arithmetic $\xi$ coefficient
intervals or universal Jensen positivity are certified by this program.
For example the sequence $(1,1,1)$ gives $J_2=(1+x)^2$ and
$H=\left(\begin{smallmatrix}2&-2\\-2&2\end{smallmatrix}\right)$,
which is PSD and singular. Adding an arbitrary positive diagonal to $H$
need not preserve $HC=C^TH$.

\section{Research tasks with mathematical acceptance criteria}
\paragraph{Track A: repair the determinant interface.}
Start with the exact diagonal Euler transfer on $\Rea s>1$.
Determine which regularized operator determinant could continue the
reciprocal Euler product without poles or zeros being silently exchanged.
Prove a relation to a self-adjoint operator before replacing the transfer
by its Hermitian part. Acceptance requires an identity of functions,
not similarity of zero plots. The displayed $2\times2$ example rejects
unrestricted determinant preservation immediately.

\paragraph{Track B: trace norms and normal families.}
Choose a connected domain containing an Euler-product overlap and the
critical strip, after removing the known poles in the target. Supply
estimates uniform in $N$ on every compact set. Either construct a
trace-class holomorphic family with uniform trace-norm convergence or
prove the hypotheses of Proposition~\ref{prop:normal}. The divergence
in Theorem~\ref{thm:tail} prevents an absolute prime-sum shortcut.
Any cancellation estimate must be independently proved without a
zero-free region equivalent to the desired conclusion.

\paragraph{Track C: positivity and canonical products.}
Compare the exact spectral quadratic form with Weil positivity and the
Li coefficients for the chosen $L$-function family. First state the
test-function space and regularization, then prove equality of the
forms. Positivity of a finite matrix or of finitely many Li coefficients
is insufficient; a uniform infinite-family statement is required.
Canonical products provide a convergence benchmark, but the preceding
converse explains the circularity risk of using zeros as input.

\paragraph{Track D: make formalization match mathematics.}
The standalone Lean companion checks exact integer and finite-matrix
obstructions. A separate mathlib layer now verifies complex Euler-factor
nonvanishing, a local boundary-convergence theorem, and holomorphic
Hermitian rigidity. It does not yet formalize the Dirichlet partial-sum
convergence or an arithmetic self-adjoint operator. Further formalization
must use the actual standard $L$-functions and explicit analytic domains. An abstract theorem of the form
``if an assumed bridge implies GRH, then GRH'' is logically valid but
would not verify the bridge. Print every theorem's axioms and review
its full type and dependencies, not merely its name.

\section{Computational examples and reproducibility}
The Python checks use integers and \texttt{fractions.Fraction}, with
no random choices, model weights, scientific datasets or floating-point
zero finding. They reproduce the Iwahori witness at $\pi=2$, the determinant
change $2\to1$, the parity obstruction, and the vanishing diagonal model.
The Lean statements verify selected finite calculations and the parity
argument in the kernel of a pinned Lean 4 toolchain. Neither executable
claims to verify the full complex-analytic arguments in this article.

Both editions share the same core source. The longer edition adds these
calculations and research interfaces; extra exposition does not strengthen
an unproved hypothesis into a theorem. The source manifest and audit
coverage record distinguish full text review, focused review, and
unverified dependencies. Publication readiness also requires independent
mathematical review, confirmed authorship, and journal-specific preparation.
